\documentclass[11pt,a4paper]{article}

\usepackage[margin=1.6cm]{geometry}
\usepackage{titlesec}
\usepackage{enumitem}
\usepackage[hidelinks]{hyperref}
\usepackage{xcolor}

\pagenumbering{gobble}
\setlength{\parindent}{0pt}

\titleformat{\section}
  {\large\bfseries}
  {}
  {0em}
  {}
  [\titlerule]

\begin{document}

% ================= HEADER =================
{\LARGE \textbf{Vinicius Abreu}} \\[4pt]
Senior Backend Engineer \& Tech Lead $|$ APIs, Microsserviços \& Sistemas Distribuídos \\[6pt]

\href{mailto:vinicius.abreu262@gmail.com}{vinicius.abreu262@gmail.com} $|$
\href{https://www.linkedin.com/in/vs-abreu/}{LinkedIn} $|$
\href{https://github.com/vinnedev}{GitHub} \\[10pt]

% ================= INTRODUÇÃO =================
\section*{Introdução}
Software Engineer com mais de 5 anos de experiência prática construindo, mantendo e evoluindo sistemas backend com foco em APIs, microsserviços, performance, confiabilidade e arquitetura limpa. Atualmente Lead Software Engineer na Celeste AI, conduzindo decisões técnicas em backend, cloud, CI/CD e sustentação de sistemas em produção. Histórico consistente em transformar requisitos complexos em soluções simples, estáveis e escaláveis. Stack principal: Node.js, TypeScript, Go, Python, PostgreSQL, Terraform, Azure DevOps e AWS.

% ================= HABILIDADES =================
\section*{Habilidades Técnicas}
\textbf{Linguagens de Programação:} TypeScript, Python, Go, Java, JavaScript \\[3pt]
\textbf{Backend \& Mensageria:} APIs RESTful, Microsserviços, Sistemas Orientados a Eventos, RabbitMQ \\[3pt]
\textbf{Resiliência \& Confiabilidade:} Estratégias de retry com exponential backoff e jitter, idempotência, circuit breakers e dead letter queues \\[3pt]
\textbf{Cloud \& DevOps:} Azure (principal), AWS, GCP, Pipelines CI/CD, Terraform, Docker \\[3pt]
\textbf{Bancos de Dados:} PostgreSQL, MySQL, MongoDB \\[3pt]
\textbf{Engenharia de Sistemas:} Sistemas Distribuídos, Design de Sistemas, Observabilidade, Escalabilidade \\[3pt]
\textbf{Linguagens Adicionais:} Delphi (sistemas legados)

% ================= EXPERIÊNCIA =================
\section*{Experiência}

\textbf{Celeste AI} — Lead Software Engineer \hfill Jul 2025 – Atual \\
São José dos Campos, Brasil (Remoto)
\vspace{-4pt}
\begin{itemize}[leftmargin=*]
  \item Liderei o design e a estabilização de uma plataforma de IA distribuída para gravação automatizada de reuniões, transcrição, geração de resumos e fluxos de notificação.
  \item Assumi a responsabilidade de um sistema instável e de alto custo, entregando uma arquitetura robusta e adequada para produção por meio da introdução de pipelines assíncronos, orquestração baseada em filas e isolamento de falhas.
  \item Projetei fluxos resilientes utilizando mensageria com estratégias de retry, exponential backoff, jitter e dead letter handling, prevenindo falhas em cascata.
  \item Melhorei significativamente a confiabilidade da plataforma e reduzi a frequência de incidentes através da implementação de observabilidade, logging estruturado e definição clara de limites entre serviços.
  \item Contribuí ativamente para decisões arquiteturais com foco em escalabilidade, tolerância a falhas e manutenibilidade a longo prazo.
\end{itemize}

\vspace{8pt}
\textbf{Celeste AI} — Software Engineer \hfill Mar 2023 – Jul 2025 \\
São José dos Campos, Brasil
\vspace{-4pt}
\begin{itemize}[leftmargin=*]
  \item Projetei e implementei um fluxo completo de assistente de reuniões, abrangendo gravação, transcrição, geração de resumos e envio de e-mails por meio de processamento assíncrono.
  \item Refatorei componentes legados para permitir escalabilidade gradual sem impacto nas cargas de trabalho existentes em produção.
  \item Resolvi problemas complexos relacionados a agendamento distribuído e execução de tarefas baseadas em tempo ao normalizar fusos horários entre containers, bancos de dados e APIs externas.
  \item Desenvolvi um serviço de sincronização de calendário integrando eventos do Google Calendar com agendadores internos para provisionar automaticamente assistentes de reunião antes de sessões ao vivo.
  \item Criei um microsserviço de otimização de custos que reduziu o uso de armazenamento em até 80\% ao extrair e padronizar faixas de áudio de arquivos de vídeo antes da transcrição.
  \item Integrei planos de assinatura, cobrança recorrente e pacotes baseados em uso, suportando modelos mensais, anuais e sob demanda.
\end{itemize}

\vspace{8pt}
\textbf{Agenzia MKT} — Software Engineer \hfill Mar 2022 – Fev 2023 \\
São José dos Campos, Brasil (Híbrido)
\vspace{-4pt}
\begin{itemize}[leftmargin=*]
  \item Atuei em uma plataforma de automação de mídias sociais orientada por IA, envolvendo geração de conteúdo, agendamento e fluxos de aprovação.
  \item Desenvolvi serviços backend e integrações para pipelines automatizados de publicação.
  \item Contribuí para o desenvolvimento de um editor voltado ao cliente, permitindo customização refinada de conteúdos gerados por IA.
  \item Desenvolvi uma aplicação mobile para suportar fluxos de aprovação e agendamento de postagens.
\end{itemize}

\vspace{8pt}
\textbf{Sistema Athos} — Desenvolvedor Full Stack \hfill Dez 2021 – Fev 2023 \\
Caraguatatuba, Brasil (Híbrido)
\vspace{-4pt}
\begin{itemize}[leftmargin=*]
  \item Desenvolvi e personalizei sistemas empresariais voltados à automação comercial.
  \item Melhorei a estabilidade, performance e usabilidade de sistemas utilizados por múltiplos clientes.
  \item Atuei em stacks de frontend (React) e backend (.NET/C\#, Node.js), além de Delphi 2010 em sistemas legados.
\end{itemize}

\vspace{8pt}
\textbf{Sistema Athos} — Suporte Técnico \hfill Dez 2020 – Dez 2021 \\
Caraguatatuba, Brasil (Híbrido)
\vspace{-4pt}
\begin{itemize}[leftmargin=*]
  \item Prestei suporte técnico para instalação, treinamento e manutenção de sistemas empresariais.
  \item Identifiquei problemas críticos em sistemas fiscais e documentei soluções permanentes posteriormente adotadas internamente.
\end{itemize}

% ================= FORMAÇÃO =================
\section*{Formação}

\textbf{MBA em Engenharia de Software} — USP / ESALQ \hfill Nov 2025 – Dez 2027 \\

\textbf{Ciência de Dados} — Pontifícia Universidade Católica de Campinas \hfill Ago 2023 – Out 2023 \\

\textbf{Bacharelado em Engenharia de Software (Incompleto)} — Estácio \hfill Ago 2022 – Ago 2025 \\

% ================= INFORMAÇÕES ADICIONAIS =================
\section*{Informações Adicionais}
Perfil com forte senso de ownership, confortável atuando em contextos ambíguos, sistemas legados e desafios técnicos de alto impacto. Atuação focada em resolução pragmática de problemas, comunicação clara e entrega consistente de sistemas confiáveis em produção. Aberto a aprendizado contínuo, colaboração e evolução técnica a partir de novos desafios e contextos.

\end{document}
